\documentclass[10pt,a4paper]{amsart}
\usepackage[margin=19mm]{geometry}
\usepackage{amsmath,amssymb,mathtools}
\usepackage[scr=rsfs,cal=euler]{mathalfa}
\usepackage{xfrac}
\usepackage[unicode,hidelinks,bookmarks=false]{hyperref}
\newcommand{\ie}{i.e.\ }
\newcommand{\cf}{cf.\ }
\newcommand{\resp}{respectively\ }
\newcommand{\Subsection}{\subsection}
\providecommand{\nobreakdash}{\nobreak\hbox{-}}
\setlength{\emergencystretch}{2em}
\multlinegap0pt
\usepackage{amssymb}
\usepackage{enumerate}
\usepackage{bm}
\newtheorem{defi}{Definition}
\newtheorem{lem}{Lemma}
\def\Cf{C_f}
\def\Ct{C_\theta}
\def\D{\mathrm{D}}
\def\E{\mathbb{E}}
\def\Id{\mathrm{Id}}
\def\R{\mathbb{R}}
\def\Rf{R_f}
\def\Rt{R_\theta}
\def\Tr{\operatorname{Tr}}
\def\cS{\mathcal{S}}
\def\d{\mathrm{d}}
\def\eps{\varepsilon}
\def\nI{(\mathrm{I})}
\def\nII{(\mathrm{II})}
\def\nIII{(\mathrm{III})}
\def\nIV{(\mathrm{IV})}
\def\nV{(\mathrm{V})}
\def\nVI{(\mathrm{VI})}
\newcommand{\norm}[1]{\|#1\|}
\newcommand{\normop}[1]{\|#1\|_{\mathrm{op}}}
\def\pCt{\Phi_{\Ct}}
\def\pRf{\Phi_{\Rf}}
\def\pRfs{\Phi_{\Rf^*}}
\def\pRt{\Phi_{\Rt}}
\newcommand{\prn}[1]{\left({#1}\right)} % parentheses
\def\trsfrmA{\mathcal{T}_{\mathcal{S}_\xi \to {\mathcal{S}}_\theta}}
\title{High-dimensional learning dynamics of multi-pass Stochastic Gradient Descent in multi-index models}
\author{Zhou Fan, Leda Wang}
\date{Source: arXiv:2601.21093v2 (CC BY 4.0)}
\begin{document}
\maketitle

\noindent\emph{Source and licence.} High-dimensional learning dynamics of multi-pass Stochastic Gradient Descent in multi-index models. Version arXiv:2601.21093v2 (CC BY 4.0), distributed by arXiv under the Creative Commons Attribution 4.0 International licence, http://creativecommons.org/licenses/by/4.0/, as recorded in the arXiv licence metadata for this exact version and shown on the abstract page https://arxiv.org/abs/2601.21093v2. Full licence notice: \texttt{licenses/arxiv-2601.21093-CC-BY-4.0.txt}.\par\smallskip

\noindent\textit{Editorial scope.} Counted unit: the article's lem lem:solution-in-space (source0.tex lines 2084--2091). The article's complete original proof of it is delivered with it (source0.tex lines 2092--2234, one \texttt{proof} environment of 143 lines). No mathematics is written, repaired or re-derived: the statement and the proof are the article's own lines, in the article's own notation, and no retained line is substituted or re-flowed.  Retained uncounted as the source-local context the proof needs: lines 713--720; lines 727--738; lines 759--771; lines 1732--1738; lines 1740--1745; lines 1751--1831.  Nothing was omitted from any retained span, so the package carries no omission record.  No retained line contains a citation, so the wrapper carries no bibliography and no bibliography entry.  No retained line contains a figure environment, an image-inclusion command or drawing code, so no image is delivered and the package declares no asset dependency.  Editorial formatting only, in addition: an AMS \texttt{amsart} preamble that restates the article's own declarations and macros used by the retained lines, namely the environments \texttt{defi}, \texttt{lem} on separate counters, together with the standard packages those lines need.  The retained environments print in this document as Definition 1, Lemma 1 in order of appearance, whereas the article prints the same items with the numbering of its own full document.  The complete native source of the article as distributed in the arXiv e-print for this exact version is bundled unchanged as \texttt{sources/193571/source0.tex}.
\begin{align}
\theta^t&=\theta^0-\int_0^t \bar\eta^r \Big({\gamma}\Gamma^r
\theta^r+g(\theta^r)+{\gamma}R_f^r(\theta^{[r]})+
{\gamma}R_f^{r,*}\theta^*\Big)\d r+\sqrt{\gamma}\,u^t,
\label{eq:def-cont-theta}\\
\xi^t&={-}\int_0^t \frac{\bar\eta^r}{\bar\kappa} R_\theta^{t,r}
f(\xi^{r-},w^*,\eps)\d z^r+w^t,\label{eq:def-cont-xi}
\end{align}

\begin{align}
r_\theta^{t,s}&=\Id_k-
\int_s^t \bar\eta^r\bigg[\Big({\gamma}\Gamma^r+\D
g(\theta^r)\Big)r_\theta^{r,s}+{\gamma}
R_f^r(r_\theta^{[r],s})\bigg]\d r \text{ for } t \geq s,
\label{eq:def-cont-deriv-t}\\
&\hspace{1in}r_\theta^{t,s}=0 \text{ for } t<s,\notag\\
r_f^{t,*}&={-}\D_\xi f(\xi^t,w^*,\eps) \int_0^t
\frac{\bar\eta^r}{\bar\kappa}
R_\theta^{t,r}r_f^{r-,*}\d z^r
+\D_{w^*} f(\xi^t,w^*,\eps),\label{eq:def-cont-deriv-xi-star}
\end{align}

\begin{align}
C_\theta^{t,s}&=\E[\theta^t \otimes \theta^s] \text{ for } t,s \in [0,\infty)
\cup \{*\}
\label{eq:def-cont-C-t}\\
R_\theta^{t,s}&=\E[r_\theta^{t,s}]
\label{eq:def-cont-R-t}\\
C_f^{t,s}&=\E\Big[\int_0^t \frac{\bar\eta^r}{\bar\kappa}f(\xi^{r-},w^*,\eps)\d z^r
\otimes \int_0^s \frac{\bar\eta^r}{\bar\kappa}f(\xi^{r-},w^*,\eps)\d z^r\Big]\label{eq:def-cont-C-h}\\
R_f^t(x^{[t]})&=\E\big[r_f^t(x^{[t]})\big]\label{eq:def-cont-R-h}\\
R_f^{t,*}&=\E[r_f^{t,*}]
\label{eq:def-cont-R-h-star}\\
\Gamma^t&=\E[\D_\xi f(\xi^t,w^*,\eps)]\label{eq:def-cont-Gamma}
\end{align}

\begin{align}
\frac{\d}{\d t}\pRt(t)&=C_0\prn{\pRt(t)+\int_0^t \pRf(t-r) \pRt(r) \d r},
\quad \text{ with } \pRt(0)=C_0,\label{eq:def-pRt}\\
\pRf(t) &= C_0\prn{\int_0^t \pRt(t-r) \pRf(r) \d r +
\pRt(t)},\label{eq:def-pRf}\\
\pRfs(t)&=C_0\left(1+\int_0^t \pRt(t-r)\pRfs(r)\d r\right).\label{eq:def-pRfs}
\end{align}

\begin{align}
\frac{\d}{\d t}\pCt(t)&=C_0\left(\pCt(t)
+\int_0^t \Phi_{R_f}(t-r)\pCt(r)\d r+\Phi_{R_f^*}(t)\right),
\text{ with } \pCt(0)=C_0(1+6\gamma).
\label{eq:def-pCt}
\end{align}

\begin{defi}\label{def:S}
Fix $T>0$ and a constant $C_0>0$ (which also defines $\pRt,\pRf,\pRfs,\pCt$).
Given a finite subset $D=\{d_1,\ldots,d_m\} \subset [0,T]$, we call
\[[0,d_1),[d_1,d_2),\ldots,[d_{m-1},d_m),[d_m,T]\]
the maximal intervals of $[0,T] \setminus D$.

Let $\mathcal{S}_\xi \equiv \mathcal{S}_\xi(T,C_0)$ be the space of all
tuples $(C_f,R_f,R_f^*,\Gamma)$ such that there exists a finite set
$D \subset [0,T]$ for which:
    \begin{enumerate}
        \item $\Cf \equiv \{\Cf^{t,s}\}_{t,s \in [0,T]}$ is a symmetric covariance kernel
on $[0,T] \otimes \R^k$, i.e.\ $\Cf^{t,s}=(\Cf^{s,t})^\top
\in \R^{k \times k}$ for all $t,s \in [0,T]$,
and for any finite subset $T_0 \subset [0,T]$,
the matrix $(\Cf^{t,s})_{t,s \in T_0} \in \R^{k|T_0| \times k|T_0|}$ is
symmetric positive-semidefinite. Furthermore, we have,
\[\Cf^{0,0}=0, \qquad
\Tr \Cf^{t,t} \leq C_0 \text{ for all } 0\leq t\leq T,\]
and for each maximal interval $I$ of $[0,T]\setminus D$,
        \begin{equation}\label{eq:Cfcontinuity}
\Tr[\Cf^{t,t}-\Cf^{t,s}-\Cf^{s,t}+\Cf^{s,s}] \leq C_0|t-s|
\text{ for all } t,s \in I.
        \end{equation}
        \item $\Rf \equiv \{\Rf^t\}_{t \in [0,T]}$ is a family of bounded linear operators
$\Rf^t:L^4([0,t],\R^k) \to \R^k$.
Fixing any process $x \in L^4([0,T],\R^k)$, the map
$t \mapsto R_f^t(x^{[t]})$ is Borel-measurable, and
        \begin{equation}\label{eq:pRf-volterra}
\norm{\Rf^{t}(x^{[t]})}^2 \leq \int_0^t \pRf(t-r) \norm{x^r}^2\d r \text{ for all }
t\in [0,T].
\end{equation}
\item $\Rf^* \equiv \{\Rf^{t,*}\}_{t \in [0,T]}$ is a Borel-measurable
process of matrices in $\R^{k \times k^*}$, where
\[\|\Rf^{t,*}\|^2 \leq \pRfs(t) \text{ for all } t \in [0,T].\]
\item $\Gamma \equiv \{\Gamma^t\}_{t \in [0,T]}$ is a Borel-measurable process
of matrices in $\R^{k \times k}$, where
\[\normop{\Gamma_t} \leq M_f \text{ for all } t \in [0,T].\]
    \end{enumerate}

Let ${\mathcal{S}}_\theta \equiv {\mathcal{S}}_\theta(T,C_0)$ be the space of
all tuples $(C_\theta,R_\theta)$ such that there exists a finite set $D \subset [0,T]$
for which:
    \begin{enumerate}
        \item $\Ct \equiv \{\Ct^{t,s}\}_{t,s \in [0,T] \cup \{*\}}$ is a symmetric
covariance kernel on $([0,T] \otimes \R^k) \times \R^{k^*}$, i.e.\
$\Ct^{t,s}=(\Ct^{s,t})^\top \in \R^{k \times k}$,
$\Ct^{t,*}=(\Ct^{*,t})^\top \in \R^{k \times k^*}$, and
$\Ct^{*,*} \in \R^{k^* \times k^*}$ for all $t,s \in [0,T]$, and
for any finite subset $T_0 \subset [0,T]$,
the matrix $(\Ct^{t,s})_{t,s \in T_0 \cup \{*\}} \in \R^{(k|T_0|+k^*) \times
(k|T_0|+k^*)}$ is symmetric positive-semidefinite. Furthermore,
we have
\[\Ct^{*,*}=\E[\theta^*\theta^{*\top}],
\quad \Ct^{0,0}=\E[\theta^0\theta^{0\top}],
\quad \Ct^{0,*}=\E[\theta^0\theta^{*\top}],\]
\[\Tr \Ct^{t,t} \leq \pCt(t) \text{ for all } 0\leq t\leq T,\]
and for each maximal interval $I$ of $[0,T]\setminus D$,
        \begin{equation}\label{eq:Cthetacontinuity}
            \Tr[\Ct^{t,t}-\Ct^{t,s}-\Ct^{s,t}+\Ct^{s,s}] \leq 
(\Phi_{C_\theta}(T)+C_0^2)|t-s|
\text{ for all } t,s \in I.
        \end{equation}
        \item $\Rt \equiv \{\Rt^{t,s}\}_{t,s \in [0,T]}$ is a Borel-measurable
process of matrices in $\R^{k \times k}$, satisfying
\[\Rt^{t,s}=0 \text{ for all $s>t$},
\qquad
\|\Rt^{t,s}\|^2 \leq \pRt(t-s) \text{ for all } s \leq t,\]
and for each maximal interval $I$ of $[0,T] \setminus D$,
\[\|\Rt^{t',s}-\Rt^{t,s}\|^2 \leq \Phi_{R_\theta}(T)|t'-t|
\text{ for all } s \in [0,T] \text{ and }
t,t' \in I \text{ with } t,t' \geq s.\]
\end{enumerate}
We set
\[\cS \equiv \cS(T,C_0)=\{(C_\theta,R_\theta,C_f,R_f,R_f^*,\Gamma):
(C_\theta,R_\theta) \in \cS_\theta,\,(C_f,R_f,R_f^*,\Gamma) \in \cS_\xi\}.\]
We denote by
$\cS^{\mathrm{cont}},\cS_\xi^{\mathrm{cont}},\cS_\theta^{\mathrm{cont}}$
these spaces where the above conditions hold with $D=\emptyset$,
i.e.\ where the continuity conditions for
$C_f,C_\theta,R_\theta$ hold on the whole interval $[0,T]$. 
\end{defi}

\begin{lem} \label{lem:solution-in-space}
Fix any $T>0$. Then for any constant $C_0>0$ large enough (depending on $T$,
$M_f$, $M_g$, $M_\eta$, $\E\|\theta^0\|^2$, $\E\|\theta^*\|^2$, and
$\|g(0)\|^2$),
$\trsfrmA$ maps $\mathcal{S}_\xi(T,C_0)$ into ${\mathcal{S}}_\theta(T,C_0)$, and
$\mathcal{S}^{\mathrm{cont}}_\xi(T,C_0)$ into
${\mathcal{S}}^{\mathrm{cont}}_\theta(T,C_0)$.
\end{lem}

\begin{proof}
Assume $(C_f,R_f,R_f^*,\Gamma) \in S_\xi(T,C_0)$, and denote
$(C_\theta,R_\theta)=\trsfrmA(C_f,R_f,R_f^*,\Gamma)$.

We note that by its definition in \eqref{eq:def-cont-C-t},
the kernel
$\{C_\theta^{t,s}\}_{t,s \in [0,T] \cup \{*\}}$ is positive-semidefinite 
with $C_\theta^{*,*},C_\theta^{0,*},C_\theta^{0,0}$ as specified in Definition
\ref{def:S}. Recall from \eqref{eq:def-cont-theta} that
\begin{align}\label{eq:bound-theta-second-moment-1}
    \theta^t&=\theta^0-\int_0^t \bar\eta^r \Big({\gamma}\Gamma^r
\theta^r+g(\theta^r)+{\gamma}R_f^{r}(\theta^{[r]})+
{\gamma}R_f^{r,*}\theta^*\Big)\d r+\sqrt{\gamma}\,u^t.
\end{align}
Applying Cauchy-Schwarz and the bound $(\int_0^t f(s)\d s)^2 \leq t\int_0^t
f(s)^2 \d s \leq T\int_0^t f(s)^2 \d s$, we have $\E\|\theta^t\|^2 \leq
6(\nI+\nII+\nIII+\nIV+\nV+\nVI)$ where
\begin{align*}
\nI&=\E\|\theta^0\|^2\\
\nII&=\E\left\|\int_0^t
\bar\eta^r\gamma\Gamma^r\theta^r \d r\right\|^2
\leq \gamma^2 M_f^2 M_\eta^2 T\int_0^t \E\norm{\theta^r}^2 \d r,\\
\nIII&=\E\left\|\int_0^t \bar\eta^r g(\theta^r) \d
r\right\|^2\leq M_\eta^2 T\int_0^t \E\|g(\theta^r)\|^2 \d r
    \leq 2M_g^2 M_\eta^2 T\int_0^t \E\norm{\theta^r}^2\d r
+2M_\eta^2 T^2\norm{g(0)}^2,\,\\
\nIV&=\E\left\|\int_0^t \bar\eta^r\gamma
R_f^{r}(\theta^{[r]}) \d r\right\|^2 \leq \gamma^2 M_\eta^2 T\int_0^t \E\norm{R_f^{r}(\theta^{[r]})}^2 \d r\\
    &\overset{(*)}{\leq} \gamma^2 M_\eta^2 T\int_0^t
\int_0^r \pRf(r-r')\E \|\theta^{r'}\|^2\d r'\,\d r,\\
\nV&=\E\left\|\int_0^t \bar\eta^r \gamma R_f^{r,*}\theta^*
\d r\right\|^2
    \leq \gamma^2 M_\eta^2 T\,\E\norm{\theta^*}^2
\int_0^t \Phi_{R_f^*}(r)\d r,\\
\nVI&=\gamma\,\E\norm{u^t}^2=\gamma \Tr \Cf^{t,t} \leq \gamma C_0.
\end{align*}
Here, we have used the property \eqref{eq:pRf-volterra} of $R_f^t(\cdot)$
to bound $(*)$.
Collecting these bounds shows, for any constant $C_0>0$ large enough,
\begin{align*}
\E\|\theta^t\|^2&\leq C_0\left(1+6\gamma+\int_0^t \left(\E\|\theta^r\|^2
+\int_0^r \Phi_{R_f}(r-r')\E\|\theta^{r'}\|^2\d r'+\Phi_{R_f^*}(r) \right)\d
r\right).
\end{align*}
The definition of $\pCt$ in \eqref{eq:def-pCt} implies
\begin{equation}\label{eq:PhiCthetaexplicit}
\pCt(t)=C_0\left(1+6\gamma+\int_0^t \left(\pCt(r)
+\int_0^r \Phi_{R_f}(r-r')\pCt(r')\d r'+\Phi_{R_f^*}(r) \right)\d r\right),
\end{equation}
so comparing with the above shows
\[\Tr C_\theta^{t,t}=\E\|\theta^t\|^2 \leq \Phi_{\Ct}(t).\]
Next, let $D$ be the discontinuity set of $(\Cf,\Rf,\Rf^*,\Gamma) \in
\cS_\xi(T,C_0)$.
For each maximal interval $I$ of $[0,T] \setminus D$
and any $s,t\in I$, by an analogous application of Cauchy-Schwarz
and this bound $\E\|\theta^t\|^2 \leq \Phi_{C_\theta}(t)$, for any constant
$C_0>0$ large enough,
\begin{align}
&\Tr[\Ct^{t,t} - \Ct^{t,s} - \Ct^{s,t} + \Ct^{s,s}]
=\E \norm{\theta^t-\theta^s}^2\notag\\
&=\E\left\|\int_s^t \bar\eta^r \Big({\gamma}\Gamma^r
\theta^r+g(\theta^r)+{\gamma}R_f^{r}(\theta^{[r]})+
{\gamma}R_f^{r,*}\theta^*\Big)\d r+\sqrt{\gamma}(u^t-u^s)\right\|^2\notag\\
&\leq C_0|t-s|\left(1+\int_s^t \E\|\theta^r\|^2\d r
+\int_s^t\int_0^r \Phi_{R_f}(r-r')\E\|\theta^{r'}\|^2\d r'\d r
+\int_s^t\Phi_{R_f^*}(r)\d r\right)+C_0\E\|u^t-u^s\|^2\notag\\
&\leq C_0|t-s|\left(1+\int_0^T \Phi_{C_\theta}(r)\d r
+\int_0^T \int_0^r \Phi_{R_f}(r-r')\Phi_{C_\theta}(r')\d r' \d r
+\int_0^T \Phi_{R_f^*}(r)\d r\right)
+C_0\E\|u^t-u^s\|^2\notag\\
&\leq (\Phi_{C_\theta}(T)+C_0^2)|t-s|\label{eq:bound-cont-Ct}
\end{align}
where the last inequality applies
$\E\norm{u^t-u^s}^2=\Tr[\Cf^{t,t}-\Cf^{t,s}-\Cf^{s,t}+\Cf^{s,s}]
\leq C_0|t-s|$ by \eqref{eq:Cfcontinuity}, and the form of
$\Phi_{\Ct}(T)$ in \eqref{eq:PhiCthetaexplicit}.
This checks all conditions of Definition \ref{def:S} for $C_\theta$.

For $R_\theta$, recall from \eqref{eq:def-cont-deriv-t} that
\begin{align}\label{eq:rthetarestate}
r_\theta^{t,s}=\Id-\int_s^t \bar\eta^r\bigg[\Big({\gamma}\Gamma^r+\D
g(\theta^r)\Big)r_\theta^{r,s}+{\gamma}
R_f^r\Big(r_\theta^{[r],s}\Big)\bigg]\d r \text{ for } t \geq s.
\end{align}
Then $\|r_\theta^{t,s}\|^2 \leq 4(\nI+\nII+\nIII+\nIV)$ where
\begin{align*}
\nI&=\|\Id\|^2=k\\
\nII&=\left\|\int_s^t \bar\eta^r\gamma\Gamma^r r_\theta^{r,s}\d
r\right\|^2
\leq \gamma^2M_f^2M_\eta^2T\int_s^t \|r_\theta^{r,s}\|^2 \d r,\\
\nIII&=\left\|\int_s^t \bar\eta^r \D g(\theta^r)r_\theta^{r,s}\d
r\right\|^2
\leq M_g^2 M_\eta^2 T\int_s^t \|r_\theta^{r,s}\|^2 \d r,\\
\nIV&=\left\|\int_s^t \bar\eta^r \gamma R_f^r\Big(r_\theta^{[r],s}\Big)\d
r\right\|^2
\leq \gamma^2 M_\eta^2 T\int_s^t
\Big\|R_f^r\Big(r_\theta^{[r],s}\Big)\Big\|^2 \d r\\
&\leq \gamma^2 M_\eta^2 T\int_s^t \int_s^r \Phi_{R_f}(r-r')
\|r_\theta^{r',s}\|^2 \d r'\d r,\\
\end{align*}
the last inequality again using \eqref{eq:pRf-volterra} and the fact that
$r_\theta^{r',s}=0$ for $r'<s$. Collecting these bounds shows
\[\|r_\theta^{t,s}\|^2
\leq C_0\left(1+\int_s^t \left(\norm{r_\theta^{r,s}}^2
+\int_s^r \Phi_{R_f}(r-r')\|r_\theta^{r',s}\|^2\d r'\right) \d r\right)\]
for any constant $C_0>0$ large enough.
The definition of $\Phi_{R_\theta}(t)$ in \eqref{eq:def-pRt} implies
\begin{equation}\label{eq:PhiRthetaexplicit}
\Phi_{R_\theta}(t-s)=C_0\left(1+\int_0^{t-s}\left(
\Phi_{R_\theta}(r)+\int_0^r \Phi_{R_f}(r-r')\Phi_{R_\theta}(r')\d
r'\right)\d r\right),
\end{equation}
so comparing with the above shows
\[\|R_\theta^{t,s}\|^2=\|\E r_\theta^{t,s}\|^2\leq \E\|r_\theta^{t,s}\|^2
\leq \Phi_{R_\theta}(t-s).\]
Then, applying this bound
$\E\|r_\theta^{t,s}\|^2 \leq \Phi_{R_\theta}(t-s)$,
we have similarly for any $C_0>0$ large enough and any $t \geq t' \geq s$ that
\begin{align}
\|R_\theta^{t,s}-R_\theta^{t',s}\|^2
&\leq \E\|r_\theta^{t,s}-r_\theta^{t',s}\|^2\notag\\
&\leq C_0|t-t'|\left(\int_{t'}^t \E\|r_\theta^{r,s}\|^2 \d r
+\int_{t'}^t \int_s^r \Phi_{R_f}(r-r')\E \|r_\theta^{r',s}\|^2 \d r' \d
r\right)\notag\\
&\leq C_0|t-t'|\left(\int_s^t \pCt(r-s)\d r
+\int_s^t \int_s^r \Phi_{R_f}(r-r')\pCt(r'-s)\d r' \d
r\right)\notag\\
&\leq |t-t'| \cdot \pRt(t-s) \leq |t-t'| \cdot \pRt(T),\label{eq:Rcontbound}
\end{align}
where the last two inequalities apply \eqref{eq:PhiRthetaexplicit} and
monotonicity of $\pRt$.
This checks all conditions of Definition \ref{def:S} for $R_\theta$, so
$(\Ct,\Rt) \in {\mathcal{S}}_\theta(T,C_0)$.

We note that \eqref{eq:Rcontbound} holds for all $t,t',s \in [0,T]$ with
$t \geq t' \geq s$, irrespective of the discontinuity set
$D$ for $(\Cf,\Rf,\Rf^*,\Gamma) \in \cS_\xi(T,C_0)$.
If $(\Cf,\Rf,\Rf^*,\Gamma)\in \mathcal{S}^{\mathrm{cont}}_\xi(T,C_0)$, then
\eqref{eq:bound-cont-Ct} also holds for all $t,s \in [0,T]$, implying by the
above arguments that the conditions of Definition \ref{def:S} for
$(C_\theta,R_\theta)$ hold with $D=\emptyset$.
Thus also $(\Ct,\Rt) \in {\mathcal{S}}_\theta^\mathrm{cont}(T,C_0)$.
\end{proof}

\end{document}
